\scanpage{220}
\mspage{75}

celle de

\[
 \alpha:\unsure{\mathfrak P|\mathcal P}(Q)\longrightarrow P.
\]

\deleted{La relation $R$ se retrouve à partir \unsure{de|des}} En \unsure{fait|fais}, il suffit de se donner

\[
 \left\{\begin{aligned}
  \beta_0 &:P\longrightarrow Q\\
  \alpha_0 &:Q\longrightarrow P
 \end{aligned}\right.\qquad
 \begin{aligned}&\text{applications}\\&\text{\uline{décroissantes}}\end{aligned}
 \tag*{(5.15)}
\]

\deleted{\unsure{q}} définies par

\[
 \left\{\begin{aligned}
 \beta_0(p)&=\beta(\{p\})\\
 \alpha_0(q)&=\alpha(\{q\}),
 \end{aligned}\right.
\]

puisqu'on en déduit $\beta,\alpha$ par

\[
 \left\{\begin{aligned}
 \beta(A)&=\mathop{\mathrm{Inf}}_{p\in A}\beta(p)\\
 \alpha(B)&=\mathop{\mathrm{Inf}}_{q\in B}\alpha(q).
 \end{aligned}\right.
 \tag*{(5.16)}
\]

La relation $R$ se déduit de la donnée de $\beta_0$, ou du choix de $\alpha_0$, par

\[
 (p,q)\in R\Longleftrightarrow q\leq\beta_0(p)\Longleftrightarrow p\leq\alpha_0(q).
 \tag*{(5.17)\unsure{\textsubscript{e}|\textsubscript{a}|.}}
\]

Cela implique donc que $\alpha_0,\beta_0$ se déterminent mutuellement, par les formules

\[
 \left\{\begin{aligned}
 \alpha_0(q)&=\deleted{\unsure{\mathrm{Inf}}}\mathrm{Sup}\,
 \underbrace{\{p\in P\mid q\leq\beta_0(p)\}}_{\{q\}^{\unsure{\cap|\frown}}}\\
 \beta_0(p)&=\mathrm{Sup}\,
 \underbrace{\{q\in Q\mid p\leq\alpha_0(q)\}}_{\{p\}^{\unsure{\cup|\smile}}}.
 \end{aligned}\right.
 \tag*{(5.18)}
\]

Reste à se demander quelles sont les applications $\alpha_0:P\unsure{-|\to}Q$, ou $\beta_0:Q\unsure{\to|-}P$, qui correspondent à une \uline{``relation dualisante''}
% archive sheet number: 220
